% chapters/ch06_sop_os.tex
% Chapter 6: SOP OS — State Label to Action Space (Constitutional Wedge)
% Source: NRMO_v5_updated.pdf, Chapter 6 (lines 6317-6456 of v5_layout.txt)

\chapter{SOP OS --- State Label to Action Space
(Constitutional Wedge)}
\label{ch:sop-os}

\nrmobox{Document status}{
\begin{tabular}{ll}
\textbf{Status}: & Historical / superseded operational snapshot \\
\textbf{Layer}:  & OPERATIONS \\
\end{tabular}\\[0.3em]
This SOP defines ``abstract state labels $\to$ action space.''\\
It does \textbf{not} govern methods of psychological state inference,
induction, or guidance.
}

\nrmobox{Normative Canon supersession notice}{
This chapter preserves an older SOP generation for provenance. Its
historical mode taxonomy \texttt{CORE / VENTURE / MISSION / SHUTDOWN}
and the associated permission table are \textbf{non-controlling} where
they conflict with the Normative Canon. Current Operational Modes are
\texttt{NORMAL / VENTURE / MISSION / SAFE}. In particular, current
\texttt{MISSION} means goal-constrained execution; the defensive
historical MISSION profile in this chapter maps to \texttt{SAFE} and/or
\texttt{MISSION-DEFENSE} according to actual behaviour. True terminal or
halt semantics remain \texttt{SHUTDOWN}. The historical table and rules
below are retained to preserve specification lineage, not to redefine
current semantics.
}

% =====================================================================
\section{Article Zero (Constitutional Wedge)}
\label{sec:sop-os-wedge}

This SOP \textbf{does not prescribe any methods} for inference,
manipulation, or induction of psychological states, cognitive
characteristics, or emotional responses.

This SOP handles only abstract state labels and the definition of
action space and mode transitions corresponding to them.

\paragraph{Purpose.}
Limit the SOP to ``how to use'' and separate the psychological system
(HST-N) as a ``map.'' This ensures that the excuse ``it's training''
cannot be used to invoke coercive SOP behaviour.

% =====================================================================
\section{Terms and Premises}
\label{sec:sop-os-terms}

\begin{description}
\item[Ruin] State in which future selectability is irreversibly lost
(state of no return).
\item[Irreversibility boundary] Boundary at which action may
transition to ``state of no return.''
\item[Degrees of freedom] Total quantity of future choices. NRMO does
\emph{not} ``constrain'' but \emph{preserves and maximises} degrees
of freedom.
\item[NRMO+Z+PP] Decision / dialogue governance that maximises
degrees of freedom on the premise of Ruin avoidance.
\end{description}

\nrmobox{Important}{
Ruin avoidance is \textbf{not a system of prohibitions}. It is a
structure to maximise action space within a range that does not
destroy future degrees of freedom.
}

% =====================================================================
\section{Input (Abstract State Labels)}
\label{sec:sop-os-input}

State labels are supplied from HST-N or other observation systems.
The SOP does not question inference methods.

\begin{table}[ht]
\centering
\small
\begin{tabularx}{\linewidth}{lXX}
\toprule
\textbf{Label} & \textbf{Definition (abstract)} & \textbf{Intuition} \\
\midrule
$\mathrm{LOAD} \in \{L_0, L_1, L_2\}$
& $L_0$ = low / $L_1$ = medium / $L_2$ = high load
& Higher $\Rightarrow$ more decision error. \\

$\mathrm{VOL} \in \{V_0, V_1, V_2\}$
& $V_0$ = stable / $V_1$ = fluctuating / $V_2$ = high
& Higher $\Rightarrow$ field more likely to collapse. \\

$\mathrm{IRREV} \in \{I_0, I_1\}$
& $I_0$ = boundary far / $I_1$ = boundary near
& Near $\Rightarrow$ higher irreversible-transition risk. \\
\bottomrule
\end{tabularx}
\caption[SOP OS state labels]{Abstract state labels consumed by the
SOP. The SOP does not infer these labels itself; HST-N supplies
them.}
\label{tab:sop-os-state-labels}
\end{table}

% =====================================================================
\section{Governance (Type ZERO Mode)}
\label{sec:sop-os-governance}

\begin{table}[ht]
\centering
\small
\begin{tabularx}{\linewidth}{lXX}
\toprule
\textbf{Mode} & \textbf{Purpose} & \textbf{Default policy} \\
\midrule
CORE
& Normal operation (maximise degrees of freedom)
& Exploration / action permitted; monitor irreversibility. \\

VENTURE
& Exploration / expansion
& Try small; expand within reversible range. \\

MISSION
& Control at boundary approach
& Organise action space; avoid irreversible direction. \\

SHUTDOWN
& Stop (safe landing)
& Temporary halt / postpone / rest / third-party check. \\
\bottomrule
\end{tabularx}
\caption[Historical SOP OS mode definitions]{Historical Type ZERO mode
definitions for this preserved SOP OS generation. Superseded for current
Operational Mode semantics by the Normative Canon.}
\label{tab:sop-os-modes}
\end{table}

\nrmobox{Historical mode principle}{
The following wording is preserved from this SOP generation: mode is not
a ``constraint'' but a \emph{gear} to preserve future degrees of freedom;
CORE / VENTURE expands degrees of freedom; MISSION / SHUTDOWN avoids
irreversibility boundaries. Current implementations must use the
canonical mode taxonomy and mapping stated in the supersession notice
above.
}

% =====================================================================
\section{Action Space and Permission Determination}
\label{sec:sop-os-action-space}

The SOP defines only ``action classification'' and ``permission
rules.'' Specific techniques and psychological operations are not
handled.

\begin{table}[ht]
\centering
\small
\begin{tabularx}{\linewidth}{lXX}
\toprule
\textbf{Category} & \textbf{Description} & \textbf{Examples (abstract)} \\
\midrule

A: Reversible, low-impact
& Reversible; degrees of freedom remain even if failed
& Confirmation, deferral, options presentation. \\

B: Reversible, medium-impact
& Reversible but may increase load / fluctuation
& Proposal, negotiation, limited trial. \\

C: Boundary approach
& Easily brings irreversibility boundary closer
& Strong demands, sudden conclusions. \\

D: Irreversible
& Irreversibly reduces degrees of freedom
& Serious destruction, severance. \\

\bottomrule
\end{tabularx}
\caption[SOP OS action categories]{Action categories and their
reversibility / impact characteristics.}
\label{tab:sop-os-categories}
\end{table}

\subsection{Minimum Permission Rules}
\label{sec:sop-os-permission-rules}

\nrmobox{Historical permission matrix}{
The permission rules below belong to the superseded mode semantics of
this SOP generation. Most importantly, ``MISSION = only A'' describes the
old defensive MISSION profile and \textbf{must not} be applied to the
current goal-constrained MISSION mode. Use current Canon-conformant
$A_{\mathrm{allowed}}$ profiles for operational permission decisions.
}

\begin{description}
\item[CORE] A, B permitted; C requires caution; D prohibited.
\item[VENTURE] A, B, C permitted within reversible scope;
D prohibited.
\item[MISSION] Only A permitted; B, C, D prohibited.
\item[SHUTDOWN] A only (confirm, defer, rest); no other actions.
\end{description}

% =====================================================================
\section{No-Exception Structure}
\label{sec:sop-os-no-exception}

\nrmobox{Constitutional constraint}{
This SOP \textbf{cannot create exceptions}.\\[0.3em]
Exceptions are requests to ``apply rules only when convenient.''
Exceptions destroy irreversibility guarantees.\\[0.3em]
NRMO+Z+PP operates without exception. If an exception feels needed,
\textbf{that itself is a MISSION / SHUTDOWN signal}.
}

\noindent
The final sentence above is preserved historical terminology; under the
current Canon, defensive narrowing must be interpreted through
\texttt{SAFE} and/or \texttt{MISSION-DEFENSE}, while true halt semantics
remain \texttt{SHUTDOWN}.

% =====================================================================
\section{Minimum Execution Checklist}
\label{sec:sop-os-checklist}

\begin{enumerate}
\item Confirm state label (LOAD / VOL / IRREV).
\item Determine current mode (CORE / VENTURE / MISSION / SHUTDOWN).
\item Check if action falls in permitted category.
\item If uncertain: treat as one category higher risk
(conservative).
\item Log: which mode, which category, why.
\end{enumerate}

\noindent
The checklist is retained as historical procedure text. Its mode-selection
step must be translated to the current canonical taxonomy before use.

% =====================================================================
\section*{Connections to Other Parts}

\begin{itemize}
\item State labels are supplied by HST-N
(Chapter~\ref{ch:hst-n}).
\item Action categories A/B/C/D are reproduced and elaborated in
Section~\ref{sec:action-space-permission-c2}.
\item The historical mode definitions correspond to those preserved in
Type ZERO (Section~\ref{sec:type-zero-modes-c2}); neither overrides the
current Normative Canon.
\item The constitutional wedge against ``training-mode exception''
is the cognitive analogue of the governance--execution separation
invariant (Invariant~\ref{inv:gov-exec-sep}).
\end{itemize}
